\documentclass{ctexbook}
\input{D:/Repository/Notebook_by_Leo_Yan/note-setup.tex}
\usepackage{wrapfig}
\usepackage{multicol}
\usepackage{paracol}
\usepackage{varwidth}
\usepackage{dsfont}
\usepackage{fix-cm}
\usepackage{array}
\usepackage{booktabs}
\usepackage{float}
\usepackage{pifont}
\usepackage{enumitem}
\usepackage{circuitikz}
\usepackage{tikz}
\usetikzlibrary{arrows.meta, positioning, shapes,calc}
\usepackage{standalone}
\usepackage{silence}
\usetikzlibrary{calc}
\allowdisplaybreaks
\raggedbottom
\usetikzlibrary{shapes,arrows,positioning,calc}
\usepackage{graphicx}
\usepackage{grffile}
\graphicspath{{D:/Repository/Notebook_by_Leo_Yan/Figures/}}
\let\cleardoublepage\clearpage
\usepackage{url}

%常数
\DeclareMathOperator{\e}{e} 
\DeclareMathOperator{\im}{i} 

% 概率与统计基础
\DeclareMathOperator{\Var}{Var}          % 方差
\DeclareMathOperator{\Cov}{Cov}          % 协方差
\DeclareMathOperator{\Corr}{Corr}        % 相关系数
\DeclareMathOperator{\E}{\mathbb{E}}     % 期望
\DeclareMathOperator{\Bias}{Bias}        % 偏倚
\DeclareMathOperator{\MSE}{MSE}          % 均方误差
\DeclareMathOperator{\RMSE}{RMSE}        % 均方根误差
\DeclareMathOperator{\SD}{SD}            % 标准差（或直接使用 \sigma）
\DeclareMathOperator{\SE}{SE}            % 标准误
\DeclareMathOperator{\Med}{Med}          % 中位数
\DeclareMathOperator{\Mode}{Mode}        % 众数
\DeclareMathOperator{\IQR}{IQR}          % 四分位距
\DeclareMathOperator{\Range}{Range}      % 极差
\DeclareMathOperator{\Exp}{Exp} 

% 常用统计缩写
\newcommand{\iid}{\text{i.i.d.}}          % 独立同分布，带点
\newcommand{\aee}{\text{a.e.}}             % 几乎处处
%\newcommand{\as}{\text{a.s.}}             % 几乎必然（almost surely）
\newcommand{\wpt}{\text{w.p.}}            % 以概率（with probability）
\newcommand{\rv}{\text{r.v.}}             % 随机变量

% 分布与特征函数
\DeclareMathOperator{\Supp}{Supp}        % 支撑集
\DeclareMathOperator{\MGF}{MGF}          % 矩母函数
\DeclareMathOperator{\CF}{CF}            % 特征函数
\DeclareMathOperator{\Logit}{Logit}      % Logit 变换
\DeclareMathOperator{\Probit}{Probit}    % Probit 变换

\title{数理统计}
\author{Leo Yan}
\date{\today}

\begin{document}
\maketitle
\tableofcontents
\clearpage
\chapter{基本概念与统计模型}

\section{基本概念}

\begin{definition}
    \textbf{总体}由与研究问题有关的所有个体组成。样本中的个体数目称为\textbf{样本大小}，抽取样本的行为称为\textbf{抽样}。

    根据总体个数有限与无限，分为有限总体和无限总体。
\end{definition}

总体也可看作所有个体上某种数量指标的集合，可以用一个随机变量及其概率分布（总体分布）描述。当数量指标有多个时，用随机向量表示总体。

\begin{definition}
    样本$\mathbf{X}=(X_1,\cdots,X_n)$可能取值的全体，构成\textbf{样本空间}，记为$\mathcal{X}$。
\end{definition}

样本的两重性：既可看做随机变量（随机向量）（抽样前，Uppercase），又可看作确定的数（抽样后，lowercase）。

简单随机抽样：(1)代表性：等概率抽选；(2)独立性：样本个体值不影响其他值。

\begin{definition}[简单随机样本]
    总体（分布）F中的容量为n的样本样本$X_1,\cdots,X_n$i.i.d.~F称为\textbf{简单随机样本}（简单样本/随机样本）。
\end{definition}

\section{统计推断}

\begin{definition}
    将出现在样本分布中的未知常数称为\textbf{参数（向量）}，如$N(a,\sigma^2)\to (a,\sigma^2)$.

    参数的取值范围称为\textbf{参数空间}。
\end{definition}

\begin{definition}
    样本分布包含未知参数时，参数取不同值的样本分布构成分部族：
    $$\mathcal{F} =\{F_{\boldsymbol{\theta}},\boldsymbol{\theta}\in\Theta\}\text{或}\mathcal{F} =\{f_{\boldsymbol{\theta}},\boldsymbol{\theta}\in\Theta\}$$
    其中$\Theta\in\mathbb{R}^k$为参数空间。
\end{definition}

\begin{example}
    指数分布$\exp(\lambda)$的样本分布族为
    $$\mathcal{F} =\{f_{\lambda}(x)=f(x,\lambda):\lambda>0\}$$
\end{example}

\begin{definition}
    从总体中抽取一定大小的样本去推断总体的概率分布的方法称为\textbf{统计推断}(statistical inference)。当样本分布形式已知，而含有未知实参数时的统计推断问题称为\textbf{参数统计推断问题}，否则称为非参数统计推断问题。
\end{definition}

参数统计推断有多种，例如参数估计和假设检验。

\section{统计量}

\begin{definition}
    由样本算出的量（样本的函数）称为\textbf{统计量}(statistic)。
\end{definition}

\begin{remark}
    \begin{itemize}
        \item 统计量只允许与样本有关而不允许与未知参数有关
        \item 统计量也具有两重性
    \end{itemize}
\end{remark}

\subsection{常见统计量}

\begin{definition}
    设$X_1,\cdots,X_n$是从某总体X中抽取的样本，称
    $$\bar{X}=\frac{1}{n}\sum_{i=1}^n X_i$$
    为\textbf{样本均值}(sample mean). 称
    $$S^2=\frac{1}{n-1}\sum_{i=1}^n (X_i-\bar{X})^2$$
    为\textbf{样本方差}(sample variance).称
    $$a_{n,k}=\frac{1}{n}\sum_{i=1}^n X_i^k,k=1,2,\cdots$$
    为\textbf{样本k阶原点矩}，$a_{n,1}=\bar{X}$。称
    $$m_{n,k}=\frac{1}{n}\sum_{i=1}^n (X_i-\bar{X})^k$$
    为\textbf{样本k阶中心距}。
\end{definition}

\begin{proposition}
    (1)$\sum_{i=1}^n (X_i-\bar{X})=0$

    (2)设$a,b$为常数，$Y_i=aX_i+b$，则
    $$\bar{Y}=a\bar{X}+b,S_Y^2=a^2 S_X^2$$

    (3)$\forall \text{常数}c$，有
    $$\sum_{i=1}^n(X_i-c)^2\ge\sum_{i=1}^n (X_i-\bar{X})^2$$
    即均值是偏差平方最小准则下的最佳估计
\end{proposition}

\begin{definition}[二维随机向量的样本矩]
    设$(X_1,Y_1),\cdots,(X_n,Y_n)$为二维样本$F(x,y)$中抽取的样本，则已定义$\bar{X},\bar{Y},S^2_X,S_Y^2$.
    $$S_{XY}=\frac{1}{n}(X_i-\bar{X})(Y_i-\bar{Y})$$
    称为X和Y的\textbf{样本协方差}(sample covariance).
    $$r=\frac{\sum_{i=1}^n (X_i-\bar{X})(Y_i-\bar{Y})}{\sqrt{[(n-1)S_X^2]\cdot [(n-1)S_Y^2]}}$$
    称为\textbf{样本相关系数}(sample covariance coefficient).
\end{definition}

\begin{definition}
    设$X_1,\cdots,X_n$是从某总体F中抽取的样本，按大小排序为$X_{(1)},\cdots,X_{(n)}$，称$(X_{(1)},\cdots,X_{(n)})$是样本$(X_1,\cdots,X_n)$的\textbf{次序统计量}(order statistics).$(X_{(1)},\cdots,X_{(n)})$的一部分也可称为次序统计量。
\end{definition}

\begin{definition}
    \begin{itemize}
        \item \textbf{样本中位数}(sample median)：$$m_{\frac{1}{2}}=\begin{cases}
        X_{((n+1)/2)},n\text{为奇数}\\
        \dfrac{1}{2}[X_{(n/2)}+X_{(n/2+1)}],n\text{为偶数}
        \end{cases}$$
        \item \textbf{样本p分位数}$(0<p<1)$(sample p-fractile)：$$X_{(m)},m=[(n+1)p]$$
        \item \textbf{样本极值}(extremum of sample)：$X_{(1)},X_{(n)}$
        \item \textbf{样本极差}(sample range)：$R=X_{(n)}-X_{(1)}$反映分散程度
    \end{itemize}
\end{definition}

\begin{definition}
    \textbf{样本变异系数}(sample coefficient of variation)与\textbf{总体变异系数}(population coefficient of variation)：$$\hat{\nu}=\frac{S_n}{\bar{Xt}},\nu=\frac{\sqrt{\mathrm{Var}(X)}}{\mathbb{E} (X)}$$
    后者以总体均值为单位度量总体分布的分散程度，前者反映后者的信息。
\end{definition}

\begin{definition}
    \textbf{样本偏度}(sample skewness):
    $$\hat{\beta}_1=\frac{m_{n,3}}{m_{n,2}^{3/2}}=\frac{\sqrt{n}\sum_{i=1}^{n}(X_i-\bar{X})^3}{\left(\sum_{i=1}^n(X_i-\bar{X})^2\right)^{3/2}}$$
    \textbf{总体偏度}(population skewness):
    $$\beta_1=\frac{\mu_3}{\mu_2^{3/2}}$$
    其中$\mu_k$为总体的k阶中心矩。

    后者衡量总体分布的非对称性，前者反映后者的信息。
\end{definition}

\begin{definition}
    \textbf{样本峰度}(sample kurtosis):
    $$\hat{\beta}_2=\frac{m_{n,4}}{m_{n,2}^{2}}-3=\frac{n\sum_{i=1}^{n}(X_i-\bar{X})^4}{\left(\sum_{i=1}^n(X_i-\bar{X})^2\right)^{2}}-3$$
    \textbf{总体峰度}(population kurtosis):
    $$\beta_2=\frac{\mu_4}{\mu_2^{2}}-3$$
    其中$\mu_k$为总体的k阶中心矩。

    后者衡量概率密度函数在众数（最大值）附近尖锐程度，前者反映后者的信息。
\end{definition}

\subsection{经验分布函数}

\begin{definition}
    设有$F(x)$的次序统计量$X_{(1)},\cdots,X_{(n)}$，称函数
    $$F_n(x)=\begin{cases}
        0,x<X_{(1)}\\
        \dfrac{k}{n},X_{(k)}\le x <X_{(k+1)},k=1,2,\cdots,n-1\\
        1,X_{(n)}\le x
    \end{cases}$$
    为\textbf{经验分布函数}(empirical distribution function).
\end{definition}

经验分布函数单调不减、右连续，是$F(x)=\mathbb{P} (X\le x)$的统计量。

定义\textbf{示性函数}
$$I_A(x)=\begin{cases}
    1,x\in A\\
    0,otherwise
\end{cases}$$
则
$$F_n(x)=\frac{1}{n}\sum_{i=1}^n I_{(-\infty,x]}(X_i)$$

在时刻$x$以前，$X_i$发生了吗？
\begin{align*}
    Y_i(x)=I_{(-\infty,x]}(X_i)&\implies P(Y_i(x)=1)=F(x),Y_1,\cdots,Y_n \text{i.i.d.}\sim B(1,F(x))\\
    &\implies nF_n(x)=\sum_{i=1}^n Y_i\sim B(n,F(x))
\end{align*}
由此可得$F_n(x)$的大样本性质：
\begin{theorem}
    (1)由中心极限定理，
    $$\frac{\sqrt{n}(F_n(x)-F(x))}{\sqrt{F(x)(1-F(x))}}\xrightarrow{\mathcal{L}}N(0,1),n\to\infty$$
    其中$\xrightarrow{\mathcal{L}}$表示依分布收敛

    (2)由Bernoulli（或辛钦）大数定律，
    $$F_n(x)\xrightarrow{\mathbb{P}}F(x),n\to\infty$$

    (3)由Borel大数定律，以概率1收敛（a.e.收敛）：
    $$\mathbb{P} (\lim_{n\to\infty}F_n(x)=F(x))=1$$

    (4)(Glivenko-Cantelli定理)用$D_n$衡量$F_n(x)$总体上靠近$F(x)$的程度：
    $$D_n=\sup_{-\infty<x<\infty}|F_n(x)-F(x)|$$
    则有
    $$\mathbb{P} (\lim_{\to\infty}D_n=0)=1$$
\end{theorem}

\section{统计模型}

一个问题的\textbf{统计模型}(statistical model)就是研究该问题时所抽样本的分布，又称概率模型，例如正态（分布）模型、泊松（分布）模型。

\section{指数型分布族}

\begin{definition}
    设$\mathcal{F}=\{f(x,\boldsymbol{\theta}):\boldsymbol{\theta}\in\Theta\}$是定义在样本空间$\mathcal{X} $上的分布族，，若概率密度函数可以表示为以下形式：
    $$f(x,\boldsymbol{\theta})=C(\boldsymbol{\theta})\exp\left[\sum_{i=1}^k Q_i(\boldsymbol{\theta})T_i(x)\right]h(x)$$
    则称此分布族为\textbf{指数型分布族/指数族}(exponential family)，其中$k\in\mathbb{N}^*,C(\boldsymbol{\theta})>0,h(x)>0$.
\end{definition}

\begin{remark}
    指数族中的所有概率密度函数$f(x,\boldsymbol{\theta})$的对$x$的支集必然不依赖于$\boldsymbol{\theta}$.
\end{remark}

\begin{definition}
    如果指数族的概率密度函数有如下形式：
    $$f(x,\boldsymbol{\theta})=C^*(\boldsymbol{\theta})\exp\left[\sum_{i=1}^k\theta_i T_i(x)\right]h(x)$$
    则称它为指数族的\textbf{自然形式}(natural form)，此时集合
    $$\Theta^*=\left\{(\theta_1,\cdots,\theta_k):\int_{\mathcal{X}}\exp\left[\sum_{i=1}^k\theta_i T_i(x)\right]h(x)\,ds<\infty\right\}$$
    称为\textbf{自然参数空间}(natural parametric space).
\end{definition}

\begin{remark}
    其中的参数组$\boldsymbol{\theta}$可能不是原分布中的参数组，而是令$\varphi_i=Q_i(\boldsymbol{\theta})$再改写符号$\boldsymbol{\varphi}$为$\boldsymbol{\theta}$.
\end{remark}

\begin{theorem}
    指数族的自然形式下，自然参数空间为凸集。
\end{theorem}

\begin{theorem}
    设指数族的自然形式下，自然参数空间有内点，记$\Theta_0=\mathrm{int}\Theta^*$.设实函数$g(x)$使得以下积分在$\Theta_0$内存在且有限：
    $$G(\boldsymbol{\theta})=\int_{\mathcal{X} }g(x)\exp\left[\sum_{j=1}^k \theta_j T_j(x)\right]h(x)\,dx$$
    则$G(\boldsymbol{\theta})$的任意偏导数存在且可以与积分号换序。
\end{theorem}

\section{贝叶斯模型}

全概率+条件概率$\implies$
\begin{theorem}[贝叶斯公式(Bayes formula)]
    设$\{B_!,\cdots,B_n\}$是样本空间$\Omega$的完备事件群（分划），$A$为$\Omega$中的事件，且$\mathbb{P}(B_i)>0,\mathbb{P} (A)>0$，则
    $$\mathbb{P}(B_i|A)=\frac{\mathbb{P} (A|B_i)\mathbb{P} (B_i)}{\mathbb{P} (A)}=\frac{\mathbb{P} (A|B_i)\mathbb{P} (B_i)}{\sum_{j=1}^n \mathbb{P}(A|B_j)(B_j)},i=1,\cdots,n$$
\end{theorem}
\begin{remark}
    $\mathbb{P} (B_i)$为A未发生时对时间$B_i$发生可能性大小的估计，而$\mathbb{P} (B_i|A)$是在已知A发生后（新信息），对$B_i$的发生可能性大小的新的认识。

    “结果”推“原因”：已知结果A，估计是哪一个原因导致A发生。
\end{remark}

统计推断使用的三种信息：\begin{itemize}
    \item 总体信息
    \item 样本信息
    \item 先验信息：抽样之前有关统计推断问题中未知参数的一些信息
\end{itemize}
基于前两者：经典统计学；基于全部三者：Bayes统计学。

\begin{definition}
    参数空间$\Theta$上$\theta$的任一概率分布称为\textbf{先验分布}(prior distribution)，记为$\pi(\theta)$.
\end{definition}

\begin{definition}
    在获得样本$\mathbf{X}=\mathbf{x}$后，$\theta$的\textbf{后验分布}(posterior distribution)记为$\pi(\theta|\mathbf{x})$，满足
    $$\pi(\theta|\mathbf{x})=\frac{f(\mathbf{x},\theta)}{m(\mathbf{x})}=\frac{f(\mathbf{x}|\theta)\pi(\theta)}{\int_{\Theta}f(\mathbf{x}|\theta)\pi(\theta)\,d\theta}$$
    其中$f(\mathbf{x}|\theta)$为样本$\mathbf{x}$的条件概率密度函数，$f(\mathbf{x},\theta)$为$\mathbf{X},\theta$的联合概率密度函数，$m(\mathbf{x})$是$\mathbf{X}$的边缘分布。离散情形下，这个公式就是Bayes公式。
\end{definition}
先验分布是抽取样本$\mathbf{X}$之前对参数$\theta$的认识，后验分布是抽取样本之后对$\theta$的认识。

Bayes学派：对参数$\theta$的统计推断必须基于$\theta$的后验分布。

先验分布的获取：\begin{itemize}
    \item 主观概率
    \item 无信息先验
    \item 共轭先验
\end{itemize}

\subsection{无信息先验}

$\Theta$有界时，无信息先验即均匀分布；$\Theta$无界时可以使用广义先验分布：

\begin{definition}
    设随机变量$X\sim f(x|\theta),\theta\in\Theta$，若$\theta$的先验密度$\pi(\theta)$满足\begin{enumerate}
        \item $\pi(\theta)\ge 0$且$\int_{\Theta}\pi(\theta)\,d\theta=\infty$
        \item 后验密度$\pi(\theta|x)$是概率密度函数
        \end{enumerate}
        则称$\pi(\theta)$为$\theta$的\textbf{广义（无信息）先验密度}(improper prior density).
\end{definition}

常取$\pi(\theta)\equiv 1$（乘以正常数无影响），此时
$$\pi(\theta|\mathbf{x})=\frac{f(\mathbf{x}|\theta)}{\int_{\Theta}f(\mathbf{x}|\theta)\,d\theta}$$

\textbf{位置参数族}：概率密度函数形如$f(x-\theta)$，$\theta$为位置参数。认为$X$与$Y=X+c$有相同的统计问题结构，于是有相同的无信息先验，此时$\pi(\theta)\equiv 1$.

\textbf{刻度参数族}：总体分布族的密度函数形如$\sigma^{-1}\varphi(x/\sigma)$，$\sigma>0$为刻度参数认为$X$与$Y=cX(c>0)$有相同的统计问题结构，于是有相同的无信息先验，此时$\pi(\sigma)=1/\sigma$.

\subsection{共轭先验}

\begin{definition}
    设$\mathcal{F}$表示又$\theta$的先验分布$\pi(\theta)$构成的分布族，如果$\forall \pi\in\mathcal{F} $及样本$\mathbf{x}$，后验分布$\pi(\theta|\mathbf{x})\in\mathcal{F}$，则称$\mathcal{F}$是一个\textbf{共轭先验分布族}(conjugate prior distribution family).
\end{definition}

共轭先验分布易于计算，并且其后验分布中一些参数可以得到很好的解释。

\chapter{抽样分布和数据的简化}

\begin{definition}
    \textbf{统计量}是样本的函数，原则上可由样本分布推出其概率分布，称之为\textbf{抽样分布}。
\end{definition}

给出抽样分布的封闭表达式，称为\textbf{精确抽样分布}，但通常难以求出，且躲在正态分布条件下得到；多数情况下研究\textbf{极限分布}，样本量大时近似精确分布。

好的统计量能够简化数据，且信息无损失，此时称之为\textbf{充分统计量}。

\section{正态总体的样本均值和样本方差的分布}

\subsection{正态变量的线性组合}

\begin{theorem}
    设$X_1,\cdots,X_n$ i.i.d.$\sim N(a_k,\sigma_k^2)$，$c_1,\cdots,c_k$为常数，记$T=\sum_{k=1}^n c_k X_k$，则$T\sim N(\mu,\tau^2)$，其中$\mu=\sum_{k=1}^n c_k a_k,\tau^2=\sum_{k=1}^n c_k^2 \sigma^2_k$.
\end{theorem}

\begin{corollary}
    设$a_1=\cdots=a_n=a,\sigma_1=\cdots=\sigma_n=\sigma$，则
    $$T\sim N\left(a\sum_{k=1}^n c_k^2,\sigma^2\sum_{k=1}^n c_k^2\right)$$
    令$c_1=\cdots=c_n=1/n$，则
    $$T=\bar{X}\sim N\left(a,\frac{\sigma^2}{n}\right)$$
\end{corollary}

\begin{theorem}
    设$X_1,\cdots,X_n$ i.i.d.$\sim N(a,\sigma^2)$，$\mathbf{Y}=A\mathbf{X}$，即
    $$\begin{bmatrix}
        Y_1\\
        \vdots\\
        Y_n
    \end{bmatrix}=\begin{bmatrix}
        a_{11} & \cdots & a_{1n}\\
        \vdots & & \vdots \\
        a_{n1} & \cdots & a_{nn}
    \end{bmatrix}\begin{bmatrix}
        X_1\\
        \vdots\\
        X_n
    \end{bmatrix}$$
    则\begin{enumerate}
        \item $\mathbf{Y}$联合高斯，且
        $$\E[Y_i]=a\sum_{k=1}^n a_{ik},\quad \Var(Y_i)=\sigma^2\sum_{k=1}^n a_{ik}^2,\quad \Cov(Y_i,Y_j)=\sigma^2\sum_{k=1}^n a_{ik}a_{jk}$$
        \item 特别地，如果$A$是单位正交阵，则$Y_1,\cdots,Y_n$ i.i.d.$\sim N(\mu_i,\sigma^2)$，其中$\mu_i=a\sum_{k=1}^n a_{ik}$.进一步，$a=0$时$Y_1,\cdots,Y_n\iid \sim N(0,\sigma^2)$
    \end{enumerate}
\end{theorem}

\begin{theorem}
    设$X_1,\cdots,X_n$ i.i.d.$\sim N(a,\sigma^2)$，则\begin{enumerate}
        \item $\bar{X}\sim N(a,\sigma^2)$
        \item $\dfrac{(n-1)S^2}{\sigma^2}\sim \chi_{n-1}^2$
        \item $\bar{X}$与$S^2$独立
    \end{enumerate}
\end{theorem}
\begin{proof}
    1.显然。对2，令单位正交阵
    $$A=\begin{bmatrix}
        \dfrac{1}{\sqrt{n}} & \cdots & \dfrac{1}{\sqrt{n}}\\
        a_{21} & \cdots & a_{2n}\\
        \vdots & & \vdots\\
        a_{n1} & \cdots & a_{nn}
    \end{bmatrix}$$
    则有$$Y_1=\frac{1}{\sqrt{n}}\sum_{k=1}^n X_k=\sqrt{n}\bar{X}$$
    $$\sum_{k=1}^n Y_k^2=\sum_{k=1}^n X_k^2$$
    因此
    $$(n-1)S^2=\sum_{k=1}^n (X_k-\bar{X})^2=\sum_{k=1}^n X_k^2-n\bar{X}^2=\sum_{k=1}^n Y_k^2-Y_1^2=\sum_{k=2}^n Y_k^2$$
    而$Y_1,\cdots,Y_n\iid\sim N(\mu_k,\sigma^2)$，其中$\mu_i=a\sum_{k=1}^n a_{ik}$，根据$A$的行向量的正交性知$\mu_k=0$，因此
    $$\frac{(n-1)S^2}{\sigma^2}=\sum_{k=2}^n \left(\frac{Y_k}{\sigma}\right)^2\sim \chi_{n-1}^2$$

    对3，$\bar{X}$只与$Y_1$有关，$S^2$只与$Y_2,\cdots,Y_n$有关，得证。
\end{proof}

\section{卡方分布，t分布和F分布}

\subsection{卡方分布}

\begin{definition}
    设$X_1,\cdots,X_n$ i.i.d.$\sim N(0,1)$，称
    $$\xi=\sum_{k=1}^n X_k^2$$
    是自由度为n的\textbf{$\chi^2$变量}，其分布称为自由度为n的\textbf{$\chi^2$分布}。
\end{definition}

\begin{theorem}
    设$\xi\sim \chi_{n}^2$，则其概率密度函数为
    $$g_n(x)=\frac{1}{2^{n/2}\Gamma(n/2)}x^{n/2-1}\e^{-x/2},x>0$$
    即$\chi_{n}^2=\Gamma(\frac{n}{2},\frac{1}{2})$.
\end{theorem}
自由度越大，$\chi_n^2$的概率密度曲线越由“偏左”趋于“对称”（极限为正态分布）；$n=1,2$单调递减，$n\ge 3$有单峰。

将$\chi_n^2$分布的上侧$\alpha(0<\alpha<1)$分位数记为$c=\chi_{n}^2(\alpha)$，它满足$\xi\sim\chi_n^2,\mathbb{P}(\xi>c)=\alpha$，通常取较小的值。

\begin{proposition}[$\chi^2$变量的性质]
    设\rv $\xi\sim\chi_n^2$，则
    \begin{enumerate}
        \item $\xi$的特征函数为$\varphi(t)=(1-2\im t)^{-n/2}$，矩母函数为$\phi(t)=(1-2t)^{-n/2}$.
        \item $\E(\xi)=n,\Var(\xi)=2n$.
        \item 设$Y\sim\Gamma(\alpha,\lambda)$，则$Z=2\lambda Y\sim\chi_{2n}^2$.
        \item 设$Z_i\iid\sim \chi_{n_i}^2$，则$\sum_{i=1}^k Z_i\sim \chi_{n_1+\cdots+n_k}^2$.
    \end{enumerate}
\end{proposition}
\begin{proof}
    利用$\Gamma(n,\lambda)$的矩母函数为$\left(\dfrac{\lambda}{\lambda-t}\right)^n$可得。
\end{proof}

\subsection{t分布}

\begin{definition}
    设$\rv X\sim N(0,1),Y\sim \chi_n^2$且两者独立，称
    $$T=\frac{X}{\sqrt{Y/n}}$$
    为自由度为n的\textbf{t变量}，其分布称为自由度为n的\textbf{t分布},记为$t_n$。
\end{definition}

\begin{theorem}
    设$\rv T\sim t_n$，则其概率密度函数为
    $$t_n(x)=\frac{\Gamma(\frac{n+1}{2})}{\Gamma(\frac{n}{2})\sqrt{n\pi}}\left(1+\frac{x^2}{n}\right)^{-\frac{n+1}{2}}$$
\end{theorem}

$t_n$类似标准正态分布$N(0,1)$，为单峰偶函数，在$x=0$处取最大值，但是峰值相对低、尾部相对粗，$n$越大越接近$N(0,1)$.

将$t_n$分布的双侧上$\alpha(0<\alpha<1)$分位数记为$c=t_n(\alpha)$，它满足$T\sim t_n,\mathbb{P}(T>c)=\alpha$。

\begin{proposition}[$t_n$变量的性质]
    设$\rv T\sim t_n$，则\begin{enumerate}
        \item $\E[T^r]$仅当$r<n$时存在，且
        $$\E[T^r]=\begin{cases}
            n^{r/2}\dfrac{\Gamma(\frac{r+1}{2})\Gamma(\frac{n-r}{2})}{\Gamma(\frac{n}{2})\Gamma(\frac{1}{2})},r\text{为偶数}\\
            0,r\text{为奇数}
        \end{cases}$$
        特别地，$n\ge 2$时,$\E[T]=0$；$n\ge 3$时，$\Var(T)=\dfrac{n}{n-2}$
        \item $t_1$就是柯西分布：$t_1(x)=\dfrac{1}{\pi(1+x^2)}$
        \item $\lim_{n\to\infty}t_n=N(0,1)$.
    \end{enumerate}
\end{proposition}

\subsection{F分布}

\begin{definition}
    设$\rv X\sim \chi_m^2,Y\sim\chi_n^2$且两者独立，则称
    $$F=\frac{X/m}{Y/n}$$
    是自由度为m和n的\textbf{F变量}，其分布函数称为自由度为m和n的\textbf{F分布}，记为$F\sim F_{m,n}$.
\end{definition}

\begin{theorem}
    $\rv Z\sim F_{m,n}$的概率密度函数为
    $$f_{m,n}(x)=B\left(\frac{m}{2},\frac{n}{2}\right)m^{\frac{m}{2}}n^{\frac{n}{2}}x^{\frac{m}{2}-1}(n+mx)^{-\frac{m+n}{2}},x>0$$
\end{theorem}

自由度m,n有序；$f_{m,n}(x)$是偏态的，m一定,n越小，偏态越严重。

将$F_{m,n}$分布的上侧$\alpha(0<\alpha<1)$分位数记为$c=F_{m,n}(\alpha)$，它满足$F\sim F_{m,n},\mathbb{P}(F>c)=\alpha$。

\begin{proposition}[F变量的性质]
    设$Z\sim F_{m,n}$，则\begin{enumerate}
        \item $1/Z\sim F_{n,m}$
        \item $\forall 0<r<\dfrac{n}{2}$，有$$\E[Z^r]=\left(\frac{n}{m}\right)^r\frac{\Gamma(\frac{m}{2}+r)\Gamma(\frac{n}{2}-r)}{\Gamma(\frac{n}{2})\Gamma(\frac{m}{2})}$$
        特别地，$$\E[Z]=\frac{n}{n-2},n>2,\quad \Var(Z)=\frac{2n^2(m+n-2)}{m(n-2)^2(n-4)},n>4$$
        \item 若$T\sim t_m$，则$T^2\sim F_{1,m}$
        \item $F_{m,n}(1-\alpha)=\dfrac{1}{F_{n,m}(\alpha)}$
    \end{enumerate}
\end{proposition}

\subsection{几个重要推论}

\begin{corollary}
    设$X_1,\cdots,X_n$ i.i.d.$\sim \Exp(\lambda)$，则
    $$2\lambda n\bar{X}=2\lambda\sum_{k=1}^n X_k\sim\chi_{2n}^2$$
\end{corollary}
\begin{proof}
    易证$2\lambda X_k\sim \chi_2^2$，利用$\chi^2$分布的性质4即证。
\end{proof}

\begin{corollary}
    设$$\mathbf{X}=(X_1,\cdots,X_n)^T\sim N_n(\boldsymbol{\mu}_{n\times 1},\Sigma_{n\times n})$$
    其中$\Cov(\mathbf{X})=\Sigma>0$（正定），则$$U=(\mathbf{X}-\boldsymbol{\mu})^T\Sigma^{-1}(\mathbf{X}-\boldsymbol{\mu})\sim \chi_n^2$$
\end{corollary}
\begin{proof}
    $\mathbf{Y}=\Sigma^{-1/2}(\mathbf{X}-\boldsymbol{\mu})\sim N_n(\mathbf{0},I)$，其中$\Sigma=O\Lambda O^T,\Sigma^{-1/2}=O\Lambda^{-1/2}O^T$，它将多元正态分布\textbf{白化}（球化）.

    由$\chi^2$分布的定义可知$U=\mathbf{Y}^T\mathbf{Y}=\sum_{k=1}^n Y_k^2\sim \chi_n^2$.
\end{proof}

\begin{corollary}
    设$X_1,\cdots,X_n\iid\sim N(a,\sigma^2)$，则
    $$T=\frac{\sqrt{n}(\bar{X}-a)}{S}\sim t_{n-1}$$
\end{corollary}
\begin{proof}
    $\bar{X}\sim N(a,\sigma^2/n)$，将其标准化：
    $$\frac{\sqrt{n}(\bar{X}-a)}{\sigma}\sim N(0,1)$$
    又知
    $$\frac{(n-1)S^2}{\sigma^2}\sim \chi_{n-1}^2,\text{即}\frac{S^2}{\sigma^2}\sim\frac{\chi_{n-1}^2}{n-1}$$
    得到
    $$T=\frac{\sqrt{n}(\bar{X}-a)/\sigma}{\sqrt{S^2/\sigma^2}}=\frac{\sqrt{n}(\bar{X}-a)}{S}\sim t_{n-1}$$
\end{proof}

\begin{corollary}
    设$X_1,\cdots,X_m\iid\sim N(a_1,\sigma^2),Y_1,\cdots,Y_n\iid\sim N(a_2,\sigma^2)$，且$X,Y$相互独立，则
    $$T_{\omega}=\frac{(\bar{X}-a_1)-(\bar{Y}-a_2)}{S_w}\cdot\sqrt{\frac{mn}{m+n}}\sim t_{n+m-2}$$
    其中
    $$(n+m-2)S_w^2=(m-1)S_X^2+(n-1)S_Y^2$$
\end{corollary}

\begin{corollary}
    设$X_1,\cdots,X_m\iid\sim N(a_1,\sigma_1^2),Y_1,\cdots,Y_n\iid\sim N(a_2,\sigma_2^2)$，且$X,Y$相互独立，则
    $$F=\frac{S_X^2}{S_Y^2}\cdot\frac{\sigma_2^2}{\sigma_1^2}\sim F_{m-1,n-1}$$
\end{corollary}

\chapter{点估计}

\chapter{区间估计}

\chapter{参数假设检验}

\end{document}